% Source: https://artofproblemsolving.com/community/c3h2871475_crossing_the_desert
% Stable identity: aops-topic-2871475-q01
Measure distance in units of $500$ miles, so the desert has width $\dfrac85$ and a full tank drives $1$ unit. Fuel burnt equals distance driven, so we minimise the total fuel used.

\textbf{Lower bound.} For $0\le x\le\dfrac85$, let $G(x)$ be the fuel burnt beyond the point $x$. All of it must be carried forward past $x$. If $G(x)>k$ for an integer $k$, the truck must cross $x$ forward at least $k+1$ times (it carries at most $1$ unit per crossing) and back at least $k$ times. So for a short stretch $[x,x+h]$ where $k<G\le k+1$, at least $(2k+1)h$ fuel is burnt there. Hence $G$ grows, moving backward from the far side, at rate at least $2k+1$ while $k<G\le k+1$, with $G\left(\dfrac85\right)=0$.
\begin{itemize}
\item $G$ reaches $1$ at the latest $1$ unit from the far side,
\item reaches $2$ at the latest a further $\dfrac13$,
\item reaches $3$ at the latest a further $\dfrac15$.
\end{itemize}
That uses $1+\dfrac13+\dfrac15=\dfrac{23}{15}$ of the width. On the remaining $\dfrac85-\dfrac{23}{15}=\dfrac1{15}$ the rate is at least $7$, adding at least $\dfrac7{15}$. So $G(0)\ge3+\dfrac7{15}=\dfrac{52}{15}$.

\textbf{Construction.} The bound is achieved by the classical plan: ferry fuel forward in $4$ loads over the first $\dfrac1{15}$ (burning $7$ times the distance), then in $3$ loads for $\dfrac15$, $2$ loads for $\dfrac13$, and finally drive the last $1$ unit on one full tank. Each stage uses exactly the rate in the bound.

So the least distance is $\dfrac{52}{15}\cdot500=\boxed{\dfrac{5200}{3}\approx1733.3}$ miles.
